% =====================================================================
\section{小横动量：运动学与 Sudakov 对数}\label{sec:tmd-kinematics}
% =====================================================================

在部分子模型中，Drell--Yan 轻子对的横动量只来自两个湮灭部分子在各自强子中的横向运动，其分布是两个横动量分布的卷积。QCD 修改了这一图像：部分子在湮灭之前可以辐射胶子，轻子对获得反冲。若辐射是一个硬胶子（$k_T\sim Q$），反冲由 $\order{\as}$ 的部分子截面描述，这是第~12~章共线因子化在 $q_T\sim Q$ 处的应用。若 $q_T\ll Q$，起作用的是大量软的与共线的辐射：每一次辐射只带走很小的横动量，而相空间给出两个对数（一个来自软，一个来自共线），所以 $\order{\as^n}$ 的截面带有多达 $2n-1$ 个 $\ln(Q^2/q_T^2)$。这些对数必须求和，最早的讨论见 Dokshitzer--Diakonov--Troian\cite{DokshitzerDiakonovTroian1978}与 Parisi--Petronzio\cite{ParisiPetronzio1979}，系统的因子化表述来自 Collins 与 Soper（背对背喷注\cite{CollinsSoper1981}）以及 Collins、Soper 与 Sterman（Drell--Yan\cite{CollinsSoperSterman1985a}）。由于反冲是对各次辐射的横动量求和，动量守恒 $\delta^{(2)}(\mathbf q_T-\sum_i\mathbf k_{iT})$ 在横向坐标空间中变成乘积，所以理论自然地在冲量参数 $\mathbf b_T$（与 $\mathbf q_T$ 共轭）中表述。

\begin{definition}[小 $q_T$ 区域]\label{def:tmd-kin}
考虑以下三类过程，强子 $A$ 沿 $+z$ 运动（光锥分量按\cref{def:prelim-lightfront}）：
\begin{enumerate}[label=(\roman*),leftmargin=2.2em,itemsep=1pt,topsep=2pt]
\item Drell--Yan $A(P_A)+B(P_B)\to\gamma^*(q)+X$，运动学按\cref{def:dy-kin}（$Q^2=q^2$，$y$ 为轻子对的快度，$x_A=\sqrt\tau\,\ee^{y}$，$x_B=\sqrt\tau\,\ee^{-y}$），并在 $\mathbf P_{A,T}=\mathbf P_{B,T}=0$ 的强子质心系中测量轻子对的横动量 $\mathbf q_T$；
\item 半单举深度非弹散射（SIDIS）$\ell+A\to\ell'+h+X$，在靶与虚光子共线的参考系中测量末态强子的横动量 $\mathbf P_{hT}$，相应的小量是 $P_{hT}/z$（$z$ 为强子的纵向动量分数）；
\item $e^+e^-\to h_1h_2X$，两个几乎背对背的强子，相应的小量是它们偏离背对背的横动量。
\end{enumerate}
记 $q_T\equiv|\mathbf q_T|$（对 SIDIS 与 $e^+e^-$ 为上述相应的横动量），\emph{小 $q_T$ 区域}（TMD 区）指 $q_T\ll Q$，幂次计数参数取 $\lambda\equiv q_T/Q$。在 $\Lambda\ll q_T\ll Q$ 时 $\order{\as}$ 的 Drell--Yan 截面具有形式
\begin{equation}\label{eq:tmd-kin-c1c0}
\frac{\dd\sigma}{\dd Q^2\,\dd y\,\dd q_T^2}=\frac{\sigma_0}{s}\,\frac{\as}{q_T^2}\Big[c_1\ln\frac{Q^2}{q_T^2}+c_0\Big]+\order{\frac{\as}{Q^2}},\qquad \sigma_0\equiv\frac{4\pi\alpha^2}{3N_cQ^2},
\end{equation}
其中的 $c_1$、$c_0$（依赖于 $x_A$、$x_B$ 与部分子分布）由此式定义。
\uses{def:dy-kin,def:prelim-lightfront}
\end{definition}

\cref{eq:tmd-kin-c1c0}的系数可以显式写出（其推导是\cref{ex:tmd-qt-asymptotics}的内容）：对 $q\bar q$ 道与 $qg$ 道，
\begin{equation}\label{eq:tmd-qt-asym}
\frac{\dd\sigma}{\dd Q^2\,\dd y\,\dd q_T^2}\bigg|_{q_T\ll Q}=\frac{\sigma_0}{s}\,\frac{\as}{2\pi}\,\frac{1}{q_T^2}\sum_qe_q^2\Big\{\Big[2A^{(1)}\ln\frac{Q^2}{q_T^2}+2B^{(1)}\Big]q(x_A)\bar q(x_B)+(P_{qq}\otimes q)(x_A)\,\bar q(x_B)+q(x_A)\,(P_{qq}\otimes\bar q)(x_B)+(P_{qg}\otimes g)(x_A)\,\bar q(x_B)+\cdots\Big\},
\end{equation}
% VF-TMD-10
其中 $q(x)$、$g(x)$ 是 $\mu=Q$ 处的 $\MSbar$ 部分子分布，$P_{ij}$ 是领头阶分裂函数（$\frac{\as}{2\pi}P$ 归一化，\cref{thm:pdf-dglap}），“$\cdots$”是交换 $A\leftrightarrow B$ 与 $q\leftrightarrow\bar q$ 的项；按 Collins--Soper--Sterman 的归一化（展开参数 $\as/\pi$，常数取 $C_1=b_0$、$C_2=1$），
\begin{equation}\label{eq:tmd-A1B1}
A^{(1)}=C_F,\qquad B^{(1)}=-\frac32\,C_F,\qquad b_0\equiv2\ee^{-\gamma_E}.
\end{equation}
% VF-TMD-7
双对数系数 $2A^{(1)}$ 来自既软又共线的胶子，$2B^{(1)}$ 来自只共线不软的胶子；$P_{qq}\otimes q$ 中 $\frac32C_F\delta(1-z)$ 的贡献与 $2B^{(1)}$ 正好相消，$q_T\to0$ 时截面中剩下的只有“软”的对数与 $(1-z)^{-1}$ 型的分布。若把 $B^{(1)}$ 取成 $0$，在\cref{ex:tmd-qt-asymptotics}所用的数值点上 $q_T/Q=10^{-3}$ 处精确截面与\cref{eq:tmd-qt-asym}之比就从 $1.000001$ 变成 $0.896$：单对数项是可以被数值区分的。\cref{eq:tmd-qt-asym}对 $q_T$ 不可积，在 $q_T\to0$ 处的发散被 $q_T=0$ 处的虚修正抵消；物理的截面在 $q_T\lesssim\Lambda$ 处由非微扰的横向运动决定。TMD 因子化的任务就是给出一个在整个 $q_T\ll Q$ 区域、对 $\as$ 的所有阶成立的公式，它在 $\Lambda\ll q_T$ 处重现\cref{eq:tmd-qt-asym}及其高阶的对数，并把 $q_T\sim\Lambda$ 的物理归入有算符定义的强子矩阵元。
